% ==============================================================================
% 01-general.tex — Раздел 1. Общее описание
% ==============================================================================

\chapter{Общее описание}
\label{ch:general}

\section{Цель создания}
\label{sec:goal}

Цель: разработка программного комплекса на основе мультиагентной архитектуры с~гибридной моделью интеграции (программный интерфейс и браузерная автоматизация) для управления цифровым присутствием субъектов \abbr{малого и среднего бизнеса}{МСБ}. OneVoice~--- платформонезависимая мультиагентная система; пользователь взаимодействует через веб-интерфейс на естественном языке, оркестратор на базе языковой модели координирует агентов.

\section{Границы проекта}
\label{sec:scope}

\textbf{Что входит в~проект:} управление бизнес-объектами, подключение каналов (VK, Telegram, Яндекс.Бизнес), контроль консистентности данных, оркестрация действий, обработка отзывов и сообщений, журналирование.

\textbf{Что не входит в~разработку:} Google Business, 2ГИС.

\section{Окружение системы}
\label{sec:environment}

Продукт~--- самостоятельная система. Компоненты: сервис программного интерфейса, сервис оркестратора, шина сообщений NATS (протокол «агент--агент»), агенты (VK, Telegram, Яндекс.Бизнес), СУБД PostgreSQL и MongoDB, кэш Redis, хранилище файлов MinIO; внешние платформы~--- программный интерфейс VK, интерфейс бота Telegram, веб-интерфейс Яндекс.Бизнес (браузерная автоматизация).

\section{Характеристики пользователей}
\label{sec:users}

Три роли пользователей~--- таблица~\ref{tab:roles}.

\begin{table}[H]
  \caption{Роли пользователей}
  \label{tab:roles}
  \small
  \begin{tabularx}{\textwidth}{|p{2.8cm}|X|}
    \hline
    \textbf{Роль} & \textbf{Функционал} \\
    \hline
    Владелец & Полный доступ: бизнес, каналы, эталонные данные, роли, подписка, журнал \\
    \hline
    Администратор & Каналы и контент: эталонные данные, синхронизация, отзывы, публикации. Без ролей и подписки \\
    \hline
    Оператор & Просмотр сообщений и отзывов, подготовка ответов, задачи. Без изменения настроек \\
    \hline
  \end{tabularx}
\end{table}

\section{Предположения и зависимости}
\label{sec:assumptions}

\textbf{Предположения:} корректные учётные данные платформ; доступность провайдера языковой модели; стабильность веб-интерфейса Яндекс.Бизнес для браузерной автоматизации.

\textbf{Зависимости:} программные интерфейсы VK и Telegram, веб-интерфейс Яндекс.Бизнес, ключ шифрования.

\section{Термины и соглашения}
\label{sec:terms}

МСБ~--- малый и средний бизнес; МАС~--- мультиагентная система; эталонные данные~--- единый источник истины по бизнес-атрибутам; консистентность~--- соответствие данных на платформах эталону; ФТ/НФТ~--- функциональные и нефункциональные требования.
